\documentclass[a4paper]{article}
\usepackage[pdftex]{graphicx}
\usepackage[utf8]{inputenc}
\usepackage{enumerate}
\usepackage{siunitx}
\sisetup{locale=DE}
\usepackage{href-ul}
\hypersetup{
	colorlinks=true,
	linkcolor=blue,
	urlcolor=blue}
\usepackage{icomma}
\usepackage{amssymb}
\usepackage{geometry}
\geometry{a4paper, top=15mm, left=15mm, right=15mm, bottom=15mm,
	headsep=10mm, footskip=12mm}


\begin{document}
	%Title
	{\bf School assignment from physics }\\
	\begin{enumerate}[1.]
		
		%Task 1
		{\bf \item In an experiment, the dependence of the current $I$ on the voltage $U$ is examined for an iron and a constantan wire.}
		
		\begin{enumerate}[a)]
			\item Create a test sketch (=circuit sketch).
			\item The following measured values result from an experiment:
			
			\hspace{15 pt}
			
			\begin{tabular}{c|c|c|c|c|c|c|c|c|c}
				& $U$ in $\SI{}{\volt}$ & 0 & 0.30 & 0.60 & 0.90 & 1.20 & 1.50 & 1.80 & 2.10 \\
				\hline
				Iron wire & $I$ in $\SI{}{\ampere}$ & 0 & 0.48 & 0.93 & 1.38 & 1.70 & 1.98 & 2.19 & 2.31 \\
				Constantan wire & $I$ in $\SI{}{\ampere}$ & 0 & 0.09 & 0.18 & 0.27 & 0.36 & 0.45 & 0.54 & 0.63
			\end{tabular}
			
			\hspace{15 pt}
			
			Evaluate the series of measurements graphically and mathematically.
			\item Interpret the course of the two conductor characteristics.
			\item Use the particle model to explain the course of the characteristic curve of the iron wire used
		\end{enumerate}
		
		%Exercise 2
		{\bf \item An ammeter with the internal resistance $R_i=\SI{70}{\ohm}$ can measure currents up to
			Measure $\SI{3,0}{\milli\ampere}$.} The measuring range should be expanded to $\SI{30}{\milli\ampere}$. Indicate how an additional resistor must be connected and calculate its value.
		
		%Task 3
		{\bf \item A current of $I=\SI{120}{\ampere}$ flows in the starter motor of a car engine during operation.} The starter battery has an internal resistance of $R_i=\SI{0.030}{\ohm}$ and a source voltage of $U=\SI{12}{\volt}$\\
		Calculate the operating voltage during the starting process as well as the resistance of the starter and its electrical power.
		
		%Task 4
		{\bf \item Two heating coils on an electric hotplate each have a resistance of
			$R=\SI{50}{\ohm}$.} With these heating coils, three levels of different heating output can be achieved using suitable circuits.
		
		\begin{enumerate}[a)]
			\item Draw a circuit diagram for the level with the smallest and the level with the highest heating output.
			\item Calculate the total resistance of the circuit and the electrical power for both stages. (voltage $\SI{230}{\volt}$).
		\end{enumerate}
		
	\end{enumerate}
	
	\fbox{
		\begin{minipage}{0.5\textwidth}
			For the solution, please \href{https://www.okuyakl.de/phys/p10elL096/le096.pdf}{click here} or scan the QR code.\\
			Additional worksheets are available at
			
			\href{http://www.okuyakl.com}{www.okuyakl.com}
		\end{minipage}
		\hfill
		\begin{minipage}{0.4\textwidth}
			\includegraphics[width=1.5 cm]{../../viecher/zwe03}
			\includegraphics[width=3 cm]{qre096}
			\includegraphics[width=2 cm]{../../viecher/afanticon1}
			
	\end{minipage}}
	
\end{document}